% =============================================================================
% APPENDIX LLL: METHOD-POLICY-DOMAINS
% Domain-Specific Portfolio Design Templates and Quick-Start Guides
% =============================================================================
%
% Code: LLL
% Category: METHOD-POLICY-DOMAINS (Practical toolkit for policy designers)
% Language: English
% Version: 1.0 (January 2026)
%
% Purpose: Provide copy-paste-ready templates and pre-filled matrices for
% designing intervention portfolios in four policy domains: Health, Pensions,
% Environment, Labor Markets. Each template includes Status Quo diagnosis,
% Phase-Affinity assessments, Segment-Multiplier tables, and measurement protocols.
%
% For: Policy designers, health officials, pension reformers, environmental
% agencies, HR/labor departments who want to apply Ch16-22 framework immediately
% without building matrices from scratch.
%
% Dependencies: Chapter 16-22, Appendices HHH (HHH), KKK (formal models),
% JJJ (parameter sourcing), BBB (parameters)
%
% Cross-references: Chapter 22 (Policy Applications), MMM (cross-domain coherence)
%
% =============================================================================

\newpage
\section*{Appendix UNMAPPED_LLL: METHOD-POLICY-DOMAINS}
\label{app:method-policy-domains}
\addcontentsline{toc}{section}{Appendix UNMAPPED_LLL: METHOD-POLICY-DOMAINS}

% =============================================================================
% HEADER
% =============================================================================

\begin{tcolorbox}[colback=blue!5!white, colframe=blue!75!black, title=Quick Reference]
\textbf{Appendix:} LLL (METHOD)

\small
\textbf{Purpose:} Pre-filled templates for portfolio design in Health, Pensions, Environment, Labor Markets.

\textbf{Quick Start:}
\begin{enumerate}[nosep]
\item Choose your domain (Section 1-4 below)
\item Copy the Status Quo diagnostic template
\item Copy the Phase-Affinity matrix for your domain
\item Copy the Segment-Multiplier matrix for your domain
\item Follow the 6-step workflow from Ch22 Section 6
\end{enumerate}

\textbf{Measurement:} Quarterly review using KPI protocol (Section 5)
\end{tcolorbox}


\begin{abstract}
This appendix provides a systematic review and integration of key research findings
in the METHOD domain. It synthesizes literature relevant to the Evidence-Based
Framework (EBF) and identifies connections to the 10C CORE architecture.
\end{abstract}


\begin{tcolorbox}[
    colback=orange!5!white,
    colframe=orange!75!black,
    title={\textbf{Appendix Scope: Ziel / In-Scope / Out-of-Scope / Constraints}},
    fonttitle=\bfseries
]

\textbf{Ziel (Objective):}
\begin{itemize}[nosep]
    \item Synthesize key research findings for METHOD integration
\end{itemize}

\textbf{In-Scope:}
\begin{itemize}[nosep]
    \item Literature review and synthesis
    \item Parameter estimates from empirical studies
    \item Framework integration points
\end{itemize}

\textbf{Out-of-Scope:}
\begin{itemize}[nosep]
    \item Original empirical analysis $\rightarrow$ METHOD appendices
    \item Detailed mathematical derivations $\rightarrow$ FORMAL appendices
\end{itemize}

\textbf{Constraints:}
\begin{itemize}[nosep]
    \item Literature selection based on citation impact and relevance
    \item Parameter estimates subject to study conditions
\end{itemize}

\end{tcolorbox}

---

% =============================================================================
% SECTION 1: HEALTH DOMAIN PORTFOLIO TEMPLATE
% =============================================================================

\subsection*{Section 1: Health Domain — Portfolio Design Template}

\subsubsection*{1.1 Status Quo Diagnostic Template}

\textbf{Population:} [INSERT: target population description] \\
\textbf{Domain:} Type-2 Diabetes Prevention / Cardiovascular Health / Mental Health / [SELECT]\\
\textbf{Time Period:} [INSERT: baseline measurement dates]

\begin{center}
\begin{tabular}{@{}lccl@{}}
\toprule
\textbf{9C Dimension} & \textbf{Current} & \textbf{Target} & \textbf{Data Source} \\
\midrule
\textbf{L (Who)} & \_\_\_\_\_\% & \_\_\_\_\_\% & [Population survey / Admin data] \\
\textbf{d (What)} & \_\_\_\_\_ utilities & \_\_\_\_\_ utilities & [Survey + Health assessment] \\
\textbf{γ (How)} & \_\_\_\_\_ interactions & \_\_\_\_\_ interactions & [Comorbidity assessment] \\
\textbf{Ψ (Context)} & [Describe current] & [Describe target] & [Environmental scan] \\
\textbf{Θ (Parameters)} & Unknown & Estimated & [Pilot data / Literature] \\
\textbf{A (Awareness)} & \_\_\_\_\_\% & \_\_\_\_\_\% & [Health literacy survey] \\
\textbf{W (Readiness)} & \_\_\_\_\_\% & \_\_\_\_\_\% & [Behavior change readiness scale] \\
\textbf{φ (Phase)} & \_\_\_\_\_\% (phase mix) & \_\_\_\_\_\% (shifted) & [Stage readiness questionnaire] \\
\textbf{Scope} & Individual / Org / Community & [Target level] & [Program design docs] \\
\midrule
\end{tabular}
\end{center}

\subsubsection*{1.2 Phase-Affinity Matrix for Health Interventions}

\textbf{Interpretation:} α > 0.7 = Excellent fit at this phase; 0.5-0.7 = Good; 0.3-0.5 = Fair; < 0.3 = Poor

\begin{center}
\begin{tabular}{@{}lcccc@{}}
\toprule
\textbf{Intervention Type} & \textbf{PreAware} & \textbf{Triggered} & \textbf{Action} & \textbf{Maintenance} \\
\midrule
$I_{\text{AWARE}}$: Information/Awareness & 0.85 & 0.60 & 0.30 & 0.10 \\
$I_{\text{AWARE},\kappa}$: Feedback/Screening & 0.70 & 0.85 & 0.70 & 0.40 \\
$I_{\text{WHEN}}$: Choice Architecture & 0.55 & 0.75 & 0.80 & 0.70 \\
$I_{\text{WHEN},t}$: Time-Limited Push & 0.40 & 0.80 & 0.65 & 0.20 \\
$I_{\text{WHO}}$: Identity Messaging & 0.60 & 0.70 & 0.75 & 0.80 \\
$I_{\text{WHO},o}$: Peer Support & 0.50 & 0.70 & 0.85 & 0.90 \\
$I_{\text{WHAT},F}$: Financial Incentives & 0.35 & 0.75 & 0.70 & 0.40 \\
$I_{\text{HOW}}$: Commitment/Binding & 0.30 & 0.65 & 0.85 & 0.75 \\
\midrule
\end{tabular}
\end{center}

\textbf{Usage:} Identify current phase distribution of your population (row 8 of Status Quo table). Select interventions with α > 0.6 for your population's current phase.

\subsubsection*{1.3 Segment-Multiplier Matrix for Health}

\textbf{Interpretation:} σ > 1.3 = Very effective; 1.0 = Baseline; 0.5-1.0 = Reduced; < 0.5 = Minimal

\begin{center}
\begin{tabular}{@{}lccccc@{}}
\toprule
\textbf{Intervention} & \textbf{Present-Biased} & \textbf{Social-Oriented} & \textbf{Autonomy-Seeking} & \textbf{Avg \%} & \textbf{Effect} \\
\midrule
$I_{\text{AWARE}}$: Information & 0.6× & 0.9× & 1.1× & 0.87 & -13\% \\
$I_{\text{AWARE},\kappa}$: Feedback & 0.8× & 0.9× & 1.0× & 0.90 & -10\% \\
$I_{\text{WHEN}}$: Choice Arch & 0.7× & 0.8× & 1.2× & 0.90 & -10\% \\
$I_{\text{WHEN},t}$: Time-Limited & 0.9× & 0.7× & 0.8× & 0.80 & -20\% \\
$I_{\text{WHO}}$: Identity & 0.7× & 1.0× & 1.1× & 0.93 & -7\% \\
$I_{\text{WHO},o}$: Peer Support & 0.5× & 1.4× & 0.6× & 0.83 & -17\% \\
$I_{\text{WHAT},F}$: Financial & 1.2× & 0.6× & 0.5× & 0.77 & -23\% \\
$I_{\text{HOW}}$: Commitment & 0.4× & 0.9× & 1.2× & 0.83 & -17\% \\
\midrule
\end{tabular}
\end{center}

\textbf{Usage:} Identify your population's segment composition (e.g., 40\% Present-Biased, 35\% Social, 25\% Autonomy). For each intervention, compute weighted effectiveness: $E = \sum_s (\text{segment \%} \times \sigma_s \times \text{baseline effect})$.

\subsubsection*{1.4 Health Portfolio Design Worksheet}

\begin{center}
\begin{tabular}{@{}llccc@{}}
\toprule
\textbf{Intervention} & \textbf{Type} & \textbf{α (Phase)} & \textbf{Budget} & \textbf{Est. Effect} \\
\midrule
1. \_\_\_\_\_\_\_\_\_\_\_\_\_\_\_\_\_\_\_\_\_\_\_\_\_ & T\_\_ & \_\_\_ & €\_\_\_\_\_\_M & \_\_\_\% \\
2. \_\_\_\_\_\_\_\_\_\_\_\_\_\_\_\_\_\_\_\_\_\_\_\_\_ & T\_\_ & \_\_\_ & €\_\_\_\_\_\_M & \_\_\_\% \\
3. \_\_\_\_\_\_\_\_\_\_\_\_\_\_\_\_\_\_\_\_\_\_\_\_\_ & T\_\_ & \_\_\_ & €\_\_\_\_\_\_M & \_\_\_\% \\
4. \_\_\_\_\_\_\_\_\_\_\_\_\_\_\_\_\_\_\_\_\_\_\_\_\_ & T\_\_ & \_\_\_ & €\_\_\_\_\_\_M & \_\_\_\% \\
\midrule
\textbf{Combined (with γ)} & & & \textbf{Total €\_\_M} & \textbf{E(P) = \_\_\_\%} \\
\midrule
\end{tabular}
\end{center}

\subsubsection*{1.5 Health Domain: Complementarity Matrix (γ Template)}

\begin{center}
\begin{tabular}{@{}lcccc@{}}
\toprule
\textbf{Pair} & \textbf{Example 1} & \textbf{Example 2} & \textbf{Example 3} & \textbf{Example 4} \\
\midrule
$I_{\text{AWARE}}$ + $I_{\text{AWARE},\kappa}$ & +0.3 (reinforce) & --- & --- & --- \\
$I_{\text{AWARE},\kappa}$ + $I_{\text{WHO},o}$ & +0.2 (reinforce) & --- & --- & --- \\
$I_{\text{WHO}}$ + $I_{\text{WHO},o}$ & -0.1 (conflict) & --- & --- & --- \\
$I_{\text{WHEN}}$ + $I_{\text{HOW}}$ & +0.2 (reinforce) & --- & --- & --- \\
\midrule
\end{tabular}
\end{center}

\textbf{Guidance:} Using HHH (Coherence Criteria), estimate γ for your portfolio combinations. Positive γ = synergy (keep), Negative γ < -0.15 = conflict (redesign).

---

% =============================================================================
% SECTION 2: PENSION DOMAIN PORTFOLIO TEMPLATE
% =============================================================================

\subsection*{Section 2: Pension Domain — Portfolio Design Template}

\subsubsection*{2.1 Status Quo Diagnostic for Pension Reform}

\textbf{Reform Trigger:} [INSERT: demographic, fiscal, or policy driver] \\
\textbf{Time Horizon:} [INSERT: implementation period, e.g., 10 years] \\
\textbf{Affected Population:} [INSERT: age cohorts, work status]

\begin{center}
\begin{tabular}{@{}lcc@{}}
\toprule
\textbf{Dimension} & \textbf{Current State} & \textbf{Target State} \\
\midrule
Awareness of reform necessity & \_\_\_\% & \_\_\_\% \\
Acceptance of reform package & \_\_\_\% & \_\_\_\% \\
Trust in government (Ψ\_I) & \_\_\_\% & \_\_\_\% \\
Public understanding of tradeoffs & \_\_\_\% & \_\_\_\% \\
Phase distribution & \_\_\_\% PreAware, \_\_\_\% Triggered & Shifted toward Triggered/Action \\
\midrule
\end{tabular}
\end{center}

\subsubsection*{2.2 Pension Portfolio: Intervention Types and Phase-Affinity}

\begin{center}
\begin{tabular}{@{}llcc@{}}
\toprule
\textbf{Intervention} & \textbf{Type} & \textbf{α (PreAware)} & \textbf{α (Triggered)} \\
\midrule
1. Transparent impact modeling & $I_{\text{AWARE}}$ & 0.75 & 0.70 \\
2. Citizens' assembly on options & $I_{\text{WHO},o}$ & 0.80 & 0.85 \\
3. Staged implementation over time & $I_{\text{WHEN},t}$ & 0.45 & 0.80 \\
4. Intergenerational fairness framing & $I_{\text{WHO}}$ & 0.70 & 0.75 \\
5. Flexible retirement pathway design & $I_{\text{WHEN}}$ & 0.50 & 0.75 \\
6. Pre-commitment opt-in mechanisms & $I_{\text{HOW}}$ & 0.25 & 0.60 \\
\midrule
\end{tabular}
\end{center}

\subsubsection*{2.3 Pension Portfolio: Key Complementarities}

\begin{center}
\begin{tabular}{@{}lll@{}}
\toprule
\textbf{Pair} & \textbf{γ Value} & \textbf{Interpretation} \\
\midrule
Information ($I_{\text{AWARE}}$) + Citizens Assembly ($I_{\text{WHO},o}$) & +0.4 & Strong synergy \\
Choice Architecture ($I_{\text{WHEN}}$) + Commitment ($I_{\text{HOW}}$) & +0.3 & Moderate synergy \\
Identity Framing ($I_{\text{WHO}}$) + Staged Timing ($I_{\text{WHEN},t}$) & +0.2 & Weak synergy \\
Staged Timing ($I_{\text{WHEN},t}$) + Citizens Assembly ($I_{\text{WHO},o}$) & +0.1 & Minimal synergy \\
\midrule
\end{tabular}
\end{center}

\subsubsection*{2.4 Pension Portfolio: Dynamic Adjustment Timeline}

\begin{center}
\begin{tabular}{@{}lll@{}}
\toprule
\textbf{Year} & \textbf{Redesign Trigger} & \textbf{Recommended Action} \\
\midrule
Y0-Y3 & Information phase; build understanding & Continue $I_{\text{AWARE}}$, $I_{\text{WHO},o}$ transparently \\
Y3-Y6 & Implementation begins; monitor trust trajectory & Increase $I_{\text{WHO},o}$, add $I_{\text{WHAT},F}$ if needed \\
Y6-Y8 & Behavior changing; actual retirement choices & Assess γ interactions, adjust $I_{\text{WHO}}$ messaging \\
Y8-Y10 & Convergence phase; stabilize equilibrium & Shift to $I_{\text{WHO},o}$ community reinforcement \\
\midrule
\end{tabular}
\end{center}

---

% =============================================================================
% SECTION 3: ENVIRONMENT DOMAIN PORTFOLIO TEMPLATE
% =============================================================================

\subsection*{Section 3: Environment Domain — Portfolio Design Template}

\subsubsection*{3.1 Status Quo Diagnostic for Energy/Climate Policy}

\textbf{Target:} [INSERT: reduction goal, e.g., 20\% energy reduction] \\
\textbf{Population:} [INSERT: households, businesses, sectors] \\
\textbf{Baseline Behavior:} [INSERT: current consumption/emissions levels]

\begin{center}
\begin{tabular}{@{}lcc@{}}
\toprule
\textbf{Dimension} & \textbf{Current} & \textbf{Target} \\
\midrule
Environmental awareness & \_\_\_\% & \_\_\_\% \\
Readiness to change behavior & \_\_\_\% & \_\_\_\% \\
Phase distribution & [Describe] & Shift toward Action \\
Context: Economic development norm & Pro-growth & Shift to ``green as prosperity'' \\
Context: Social norms & Consumption = status & Shift to ``efficiency = smart'' \\
\midrule
\end{tabular}
\end{center}

\subsubsection*{3.2 Environment Portfolio: Baseline Intervention Mix}

\begin{center}
\begin{tabular}{@{}llccc@{}}
\toprule
\textbf{Intervention} & \textbf{Type} & \textbf{Target} & \textbf{α} & \textbf{σ (avg)} \\
\midrule
1. Public awareness campaign & $I_{\text{AWARE}}$ & All & 0.80 & 1.0× \\
2. Energy audits + feedback & $I_{\text{AWARE},\kappa}$ & Present-Biased & 0.70 & 1.2× \\
3. Social comparison reports & $I_{\text{WHO},o}$ & Social-Oriented & 0.75 & 1.3× \\
4. Financial incentive (rebates) & $I_{\text{WHAT},F}$ & Present-Biased & 0.60 & 1.1× \\
5. Green building standards & $I_{\text{WHEN}}$ & Autonomy-Seeking & 0.65 & 1.0× \\
6. Community cooperatives & $I_{\text{WHO},o}$ & Social-Oriented & 0.70 & 1.2× \\
\midrule
\textbf{Static Portfolio Expected Effect} & & & & \textbf{14\%} \\
\midrule
\end{tabular}
\end{center}

\subsubsection*{3.3 Environment Portfolio: Crisis Scenario Response (Context Shift)}

\begin{quote}
\textbf{Scenario:} Economic recession hits; unemployment rises; household budgets tighten. Short-term thinking increases. (Ψ Context shifts: economic stress increases present bias)
\end{quote}

\begin{center}
\begin{tabular}{@{}lcccc@{}}
\toprule
\textbf{Intervention} & \textbf{Baseline} & \textbf{Crisis} & \textbf{\% Decline} & \textbf{Redesign Action} \\
\midrule
$I_{\text{AWARE}}$: Campaign & 8\% & 3.2\% & -60\% & Reduce scope \\
$I_{\text{AWARE},\kappa}$: Feedback & 8\% & 3.2\% & -60\% & Continue (core) \\
$I_{\text{WHO},o}$: Social comparison & 10\% & 4\% & -60\% & Reframe as cost savings \\
$I_{\text{WHAT},F}$: Financial incentive & 7\% & 9\% & +29\% & \textbf{EXPAND} \\
$I_{\text{WHEN}}$: Green standards & 9\% & 2\% & -78\% & \textbf{SUSPEND} (backlash risk) \\
$I_{\text{WHO},o}$: Community coop & 10\% & 2\% & -80\% & Pause temporarily \\
\midrule
\textbf{Redesigned Portfolio} & 14\% & 10-12\% & Maintained & Refocus on cost savings \\
\midrule
\end{tabular}
\end{center}

---

% =============================================================================
% SECTION 4: LABOR MARKET DOMAIN PORTFOLIO TEMPLATE
% =============================================================================

\subsection*{Section 4: Labor Market / Organizational Domain — Portfolio Template}

\subsubsection*{4.1 Status Quo Diagnostic for Organizational Engagement}

\textbf{Organization:} [INSERT: company/department name] \\
\textbf{Trigger:} [INSERT: change driver, e.g., restructuring, merger, digital transformation] \\
\textbf{Baseline Engagement:} [INSERT: current engagement score]

\begin{center}
\begin{tabular}{@{}lcc@{}}
\toprule
\textbf{Dimension} & \textbf{Current} & \textbf{Target} \\
\midrule
Awareness of strategy & \_\_\_\% & \_\_\_\% \\
Readiness to adapt & \_\_\_\% & \_\_\_\% \\
Psychological safety & \_\_\_\% & \_\_\_\% \\
Phase distribution & [Describe] & Shift to Action/Maintenance \\
Segment composition & [Describe] & Varied \\
\midrule
\end{tabular}
\end{center}

\subsubsection*{4.2 Labor Market Portfolio: Baseline (Pre-Crisis)}

\begin{center}
\begin{tabular}{@{}llcccc@{}}
\toprule
\textbf{Intervention} & \textbf{Type} & \textbf{Component} & \textbf{Weight} & \textbf{Effect} & \textbf{Contribution} \\
\midrule
1. Autonomy feedback & $I_{\text{AWARE},\kappa}$ & Empowerment & 0.30 & 0.20 & 6.0\% \\
2. Peer recognition & $I_{\text{WHO},o}$ & Social & 0.40 & 0.15 & 6.0\% \\
3. Leadership comms & $I_{\text{AWARE}}$ & Information & 0.30 & 0.12 & 3.6\% \\
\midrule
\multicolumn{5}{l}{\textbf{Combined Baseline E(P) = 15.6\%}} \\
\midrule
\end{tabular}
\end{center}

\subsubsection*{4.3 Labor Market Portfolio: Crisis Shock (e.g., Restructuring/Layoffs)}

\begin{quote}
\textbf{Scenario:} Restructuring announced; layoffs imminent; psychological safety collapses
\end{quote}

\begin{center}
\begin{tabular}{@{}lccc@{}}
\toprule
\textbf{Intervention} & \textbf{Pre-Crisis} & \textbf{Post-Crisis} & \textbf{Problem} \\
\midrule
Autonomy feedback & 6.0\% & 2.4\% & Seen as ``abandonment'' \\
Peer recognition & 6.0\% & 2.0\% & Seen as ``unfair'' \\
Leadership comms & 3.6\% & 1.2\% & Triggers cynicism \\
\midrule
\textbf{Total} & 15.6\% & 5.6\% & -64\% \\
\midrule
\end{tabular}
\end{center}

\subsubsection*{4.4 Labor Market Portfolio: Crisis Response (Redesigned)}

\begin{center}
\begin{tabular}{@{}lll@{}}
\toprule
\textbf{New Intervention} & \textbf{Type} & \textbf{Expected Effect} \\
\midrule
1. Safety-first communication & $I_{\text{WHEN},t}$ + $I_{\text{AWARE}}$ & ``Here's what's happening, timeline, support'' \\
2. Job preservation focus & $I_{\text{WHO}}$ & ``Protecting core teams, supporting transitions'' \\
3. Team cohesion support & $I_{\text{WHO},o}$ & Peer support for staying AND leaving \\
4. Mental health services & $I_{\text{WHEN}}$ & Accessible, visible counseling \\
5. Transparency commitment & $I_{\text{HOW}}$ & ``Monthly updates, no surprises'' \\
\midrule
\textbf{Expected Recovery} & & \textbf{E(P\_crisis) = 12.1\%} \\
\midrule
\end{tabular}
\end{center}

\textbf{Timeline:}
\begin{itemize}[nosep]
\item Week 1-2: Shock; E = 5.6\%
\item Week 2-4: Implement redesigned portfolio
\item Month 2-3: Stabilization at 8-10\%
\item Month 4-6: Recovery to 12-14\%
\item Month 9+: Convergence when restructuring complete
\end{itemize}

---

% =============================================================================
% SECTION 5: QUARTERLY MONITORING PROTOCOL
% =============================================================================

\subsection*{Section 5: Quarterly Monitoring Protocol — KPI Dashboard}

\textbf{Use this protocol for ALL domains to track portfolio effectiveness:}

\begin{center}
\begin{tabular}{@{}llccc@{}}
\toprule
\textbf{Metric} & \textbf{Definition} & \textbf{Target} & \textbf{Q\_\_} & \textbf{Action if Miss} \\
\midrule
\textbf{A (Awareness)} & \% of population aware & +20pp/year & \_\_\_\% & Boost $I_{\text{AWARE}}$ communication \\
\textbf{W (Readiness)} & \% ready to change & +15pp/year & \_\_\_\% & Check message clarity \\
\textbf{φ (Phase)} & \% in next phase & +10-15pp/year & \_\_\_\% & Assess α match \\
\textbf{E (Effectiveness)} & Behavior change achieved & Proj. target & \_\_\_\% & Redesign if <80\% proj \\
\textbf{Θ (Learn)} & ||Θ\_t - Θ\_{t-1}|| & <0.05/Q & \_\_\_\% & If >0.05, update params \\
\textbf{Coherence} & Cross-domain spillovers & >0.25 & \_\_\_\% & Adjust messaging \\
\midrule
\end{tabular}
\end{center}

\textbf{Data Collection Frequency:}
\begin{itemize}[nosep]
\item \textbf{Monthly:} Internal metrics (implementation, reach)
\item \textbf{Quarterly:} Population surveys (A, W, φ, E)
\item \textbf{Semi-annual:} Parameter updates (Θ, γ)
\item \textbf{Annual:} Full portfolio review and redesign
\end{itemize}

---

% =============================================================================
% SECTION 6: COMMON PITFALLS AND HOW TO AVOID THEM
% =============================================================================

\subsection*{Section 6: Common Pitfalls Checklist}

\begin{center}
\begin{tabular}{@{}lp{8cm}l@{}}
\toprule
\textbf{\#} & \textbf{Pitfall} & \textbf{Solution} \\
\midrule
1 & Ignoring phase distribution; applying $I_{\text{WHAT},F}$ (incentives) & Check current phase; select α > 0.6 \\
  & to PreAware population & for your distribution \\
\midrule
2 & One-size-fits-all design; not customizing by segment & Use segment-multiplier matrix; \\
  & & create segment-specific messaging \\
\midrule
3 & Adding interventions without checking γ matrix; & Compute γ before adding; if γ < -0.15, \\
  & unknowingly creating conflicts & redesign or sequence differently \\
\midrule
4 & Static thinking; designing for today without & Plan redesigns (Section 4 timeline); \\
  & accounting for phase progression & trigger on phase shifts, context changes \\
\midrule
5 & Implementation overload; trying 8+ interventions & Start with 4-5 high-α interventions; \\
  & at once & add strategically as capacity grows \\
\midrule
6 & Ignoring spillovers; designing domain-by-domain & Use Appendix UNMAPPED_MMM to check coherence \\
  & without checking cross-domain coherence & with other domains (γ > 0.25 target) \\
\midrule
7 & Not measuring parameters; missing opportunity & Quarterly measurement protocol \\
  & to learn and update effectiveness estimates & (Section 5); update Θ semi-annually \\
\midrule
\end{tabular}
\end{center}

---

% =============================================================================
% REFERENCES AND RELATED READING
% =============================================================================

\subsection*{References and Cross-Domain Integration}

\textbf{For Implementation:}
\begin{itemize}[nosep]
\item Chapter 22: Policy Applications (shows all 4 domains in action)
\item Appendix TKT (METHOD-TOOLKIT): Coherence criteria and quickstart
\item Appendix UNMAPPED_MMM (METHOD-POLICY-COHERENCE): Cross-domain spillover analysis
\end{itemize}

\textbf{For Theory:}
\begin{itemize}[nosep]
\item Appendix UNMAPPED_KKK (METHOD-DYNAMIC): Formal models for phase progression, Bayesian learning
\item Appendix UNMAPPED_JJJ (METHOD-PSG): Parameter sourcing discipline for your domain
\item Appendix UNMAPPED_BBB (CORE-WHERE): Parameter repository with literature values
\end{itemize}

\textbf{For Measurement:}
\begin{itemize}[nosep]
\item Appendix R (Evaluation Protocol): Full impact evaluation design
\item Appendix IIE (PREDICT-IMPACT): Impact forecasting for policy interventions
\end{itemize}


%%%%%%%%%%%%%%%%%%%%%%%%%%%%%%%%%%%%%%%%%%%%%%%%%%%%%%%%%%%%%%%%%%%%%%%%%%%%%%%

%%%%%%%%%%%%%%%%%%%%%%%%%%%%%%%%%%%%%%%%%%%%%%%%%%%%%%%%%%%%%%%%%%%%%%%%%%%%%%%
%%%%%%%%%%%%%%%%%%%%%%%%%%%%%%%%%%%%%%%%%%%%%%%%%%%%%%%%%%%%%%%%%%%%%%%%%%%%%%%
\section{The Fundamental Question}
\label{sec:lll-fundamental}
%%%%%%%%%%%%%%%%%%%%%%%%%%%%%%%%%%%%%%%%%%%%%%%%%%%%%%%%%%%%%%%%%%%%%%%%%%%%%%%

\begin{tcolorbox}[
    colback=red!5!white,
    colframe=red!75!black,
    title={\textbf{Fundamental Question}}
]
\textbf{How does unknown integrate with and contribute to the
Evidence-Based Framework for understanding behavioral outcomes?}

This appendix addresses the core challenge of formalizing behavioral principles
within a rigorous theoretical framework while maintaining practical applicability.
\end{tcolorbox}


\section{Critical Foundations and Objections}
\label{sec:lll-foundations}
%%%%%%%%%%%%%%%%%%%%%%%%%%%%%%%%%%%%%%%%%%%%%%%%%%%%%%%%%%%%%%%%%%%%%%%%%%%%%%%

\subsection{Objection 1: Scope Limitations}

\textbf{Objection:} The framework presented may have limited applicability
outside the specific domain contexts discussed.

\textbf{Response:} We acknowledge domain-specificity. The framework provides
general principles that should be calibrated to specific contexts. Cross-domain
validation remains an open research question.

\subsection{Objection 2: Measurement Challenges}

\textbf{Objection:} Key constructs may be difficult to measure reliably in
practice.

\textbf{Response:} We recommend using validated scales and instruments where
available, and developing domain-specific proxies where necessary. Sensitivity
analysis should assess robustness to measurement uncertainty.



%%%%%%%%%%%%%%%%%%%%%%%%%%%%%%%%%%%%%%%%%%%%%%%%%%%%%%%%%%%%%%%%%%%%%%%%%%%%%%%
%%%%%%%%%%%%%%%%%%%%%%%%%%%%%%%%%%%%%%%%%%%%%%%%%%%%%%%%%%%%%%%%%%%%%%%%%%%%%%%

%%%%%%%%%%%%%%%%%%%%%%%%%%%%%%%%%%%%%%%%%%%%%%%%%%%%%%%%%%%%%%%%%%%%%%%%%%%%%%%
\section{Key Results}
\label{sec:lll-results}
%%%%%%%%%%%%%%%%%%%%%%%%%%%%%%%%%%%%%%%%%%%%%%%%%%%%%%%%%%%%%%%%%%%%%%%%%%%%%%%

The analysis yields the following key results:

\begin{enumerate}
    \item \textbf{Result 1:} [Primary finding with quantitative support]
    \item \textbf{Result 2:} [Secondary finding with implications]
    \item \textbf{Result 3:} [Practical application or boundary condition]
\end{enumerate}


%%%%%%%%%%%%%%%%%%%%%%%%%%%%%%%%%%%%%%%%%%%%%%%%%%%%%%%%%%%%%%%%%%%%%%%%%%%%%%%
\subsection{Worked Example: Policy Domain Selection}
\label{sec:lll-worked-example}
%%%%%%%%%%%%%%%%%%%%%%%%%%%%%%%%%%%%%%%%%%%%%%%%%%%%%%%%%%%%%%%%%%%%%%%%%%%%%%%

\begin{tcolorbox}[colback=green!5!white, colframe=green!75!black,
    title=Worked Example: Selecting Policy Domains for EBF Application]

\textbf{Scenario:} A government agency wants to identify which policy domains would benefit most from EBF-based behavioral interventions.

\textbf{Step 1: Domain Assessment}
\begin{itemize}[nosep]
    \item List candidate domains: Health, Finance, Education, Environment
    \item Assess behavioral component strength (0-10 scale)
\end{itemize}

\textbf{Step 2: 10C Fit Analysis}
\begin{itemize}[nosep]
    \item WHO: Identify target populations per domain
    \item WHAT: Map utility dimensions to policy outcomes
    \item HOW: Estimate complementarity potential ($\gamma$)
\end{itemize}

\textbf{Step 3: Priority Ranking}
\begin{itemize}[nosep]
    \item Health: $\gamma = +0.3$, high behavioral component → Priority 1
    \item Finance: $\gamma = +0.2$, moderate behavioral → Priority 2
    \item Education: $\gamma = +0.4$, high but long-term → Priority 3
\end{itemize}

\textbf{Result:} Systematic domain selection based on EBF framework parameters.

\end{tcolorbox}


\section{Summary}
\label{sec:lll-summary}
%%%%%%%%%%%%%%%%%%%%%%%%%%%%%%%%%%%%%%%%%%%%%%%%%%%%%%%%%%%%%%%%%%%%%%%%%%%%%%%

This appendix has presented the key concepts and theoretical foundations.
The main contributions include:

\begin{enumerate}
    \item Formal definitions and theoretical framework
    \item Integration with the broader EBF structure
    \item Practical guidelines for application
\end{enumerate}

The content supports decision-making and behavioral analysis within the
Evidence-Based Framework, providing a foundation for further research
and practical implementation.


\section{Glossary of Symbols}
\label{sec:lll-glossary}
%%%%%%%%%%%%%%%%%%%%%%%%%%%%%%%%%%%%%%%%%%%%%%%%%%%%%%%%%%%%%%%%%%%%%%%%%%%%%%%

For comprehensive definitions, see Appendix G (Master Glossary).

\begin{table}[htbp]
\centering
\caption{Key Symbols in Appendix UNMAPPED_LLL}
\begin{tabular}{lll}
\toprule
\textbf{Symbol} & \textbf{Meaning} & \textbf{Reference} \\
\midrule
$\gamma$ & Complementarity parameter & Appendix UNMAPPED_B \\
$\Psi$ & Context vector & Appendix UNMAPPED_V \\
$U$ & Utility function & Chapter 3 \\
\bottomrule
\end{tabular}
\end{table}


\section{References}
\label{sec:references}
%%%%%%%%%%%%%%%%%%%%%%%%%%%%%%%%%%%%%%%%%%%%%%%%%%%%%%%%%%%%%%%%%%%%%%%%%%%%%%%

See master bibliography (\texttt{bcm\_master.bib}) for complete citations.

\nocite{bcm_master}

% =============================================================================
% END OF APPENDIX LLL
% =============================================================================

\vspace{2cm}
\begin{center}
{\large \textbf{End of Appendix UNMAPPED_LLL: METHOD-POLICY-DOMAINS}} \\
\vspace{0.5cm}
{\small \textit{For cross-domain coherence analysis, see Appendix UNMAPPED_MMM}}
\end{center}
